../cictikz/src/cictikz/data/tex/cictikz_preamble.tex

% One-point modulation: the modulation f_mod is added to the reference
% before the loop, so the loop has to track it -- which only works for
% modulation well inside the loop bandwidth.  Geometry matches
% l10_freq_fb; l08_pll_2mod is the two-point version.

\begin{circuitikz}[american, thick, transform shape]

  %%%%%%%%%%%%%%%%%%%%%%%%%%%%%%%%%%%%%%%%%%%%%%%%%%%%%%%%%%%%%%%%%%%%%%
%% Shared pieces for the l04_afe signal flow graphs, l4_first_order and
%% l4_biquad.  The biquad is meant to read as the first-order graph with a
%% second integrator added, so the summing node and the 1/s block have to be
%% the same size and shape in both.  They are the same size in the original
%% artwork too.
%%
%% Names are prefixed and letters only, and computed coordinates are braced:
%% TikZ scans a coordinate for its closing parenthesis, so a '(' inside an
%% expression ends it early.
%%%%%%%%%%%%%%%%%%%%%%%%%%%%%%%%%%%%%%%%%%%%%%%%%%%%%%%%%%%%%%%%%%%%%%

\newcommand{\sfgR}{0.25}    % summing node radius
\newcommand{\sfgBw}{1.0}    % half width of the 1/s block
\newcommand{\sfgBh}{0.65}   % half height of the 1/s block

% A summing node at (#1,#2).
\newcommand{\sfgSum}[2]{%
  \draw (#1,#2) circle (\sfgR);
  \draw ({#1-0.17},#2) -- ({#1+0.17},#2);
  \draw (#1,{#2-0.17}) -- (#1,{#2+0.17});
}

% A block centred at (#1,#2) carrying #3.
\newcommand{\sfgBox}[3]{%
  \draw ({#1-\sfgBw},{#2-\sfgBh}) rectangle ({#1+\sfgBw},{#2+\sfgBh});
  \node at (#1,#2) {#3};
}

\newcommand{\sfgDot}[2]{\fill (#1,#2) circle (0.07);}


  \newcommand{\pllBox}[3]{%
    \draw ({#1-0.55},{#2-0.55}) rectangle ({#1+0.55},{#2+0.55});
    \node at (#1,#2) {#3};}
  \newcommand{\pllOsc}[2]{%
    \pllBox{#1}{#2}{}
    \draw ({#1-0.36},#2) sin ({#1-0.18},{#2+0.22}) cos (#1,#2)
          sin ({#1+0.18},{#2-0.22}) cos ({#1+0.36},#2);}

  \def\ym{0}      % the forward path
  \def\yt{1.8}    % the feedback path
  \def\xa{1.6}    % modulation sum
  \def\xs{3.3}    % loop sum
  \def\xh{5.2}    % H(s)
  \def\xo{7.2}    % oscillator
  \def\xn{6.2}    % divider on the feedback path
  \def\xt{8.7}    % where the feedback taps the output

  % reference plus modulation
  \node[anchor=east] at (0.1,\ym) {$f_{in}$};
  \draw[->] (0.1,\ym) -- ({\xa-\sfgR},\ym);
  \sfgSum{\xa}{\ym}
  \draw[->] (\xa,{\ym+1.5}) -- (\xa,{\ym+\sfgR});
  \node[anchor=south west] at ({\xa+0.1},{\ym+1.4}) {$f_{mod}$};

  % the loop
  \draw[->] ({\xa+\sfgR},\ym) -- ({\xs-\sfgR},\ym);
  \sfgSum{\xs}{\ym}
  \draw[->] ({\xs+\sfgR},\ym) -- ({\xh-0.55},\ym);
  \pllBox{\xh}{\ym}{$H(s)$}
  \draw ({\xh+0.55},\ym) -- ({\xo-0.55},\ym);
  \pllOsc{\xo}{\ym}
  \draw ({\xo+0.55},\ym) -- (10.1,\ym);
  \node[anchor=south] at (9.7,0.25) {$f_{o}$};

  % feedback through the divider
  \sfgDot{\xt}{\ym}
  \draw (\xt,\ym) -- (\xt,\yt) -- ({\xn+0.55},\yt);
  \pllBox{\xn}{\yt}{$\div N$}
  \draw ({\xn-0.55},\yt) -- (\xs,\yt);
  \draw[->] (\xs,\yt) -- (\xs,{\ym+\sfgR});
  \node[anchor=south west] at ({\xs+0.12},{\ym+0.32}) {$-$};

\end{circuitikz}
\end{document}
